

\documentclass[12pt]{article}

\usepackage[margin = .8in]{geometry}
\usepackage{amsmath}
\usepackage{graphicx}
\usepackage{multicol, enumerate, tabularx}

\usepackage[parfill]{parskip}

\usepackage{adjustbox, soul}

\usepackage{fancyhdr}
\pagestyle{fancy}

\lhead{\emph{Math F113X Cryptograph 4: Double Transposition}}
\rhead{\emph{lecture notes}}

\renewcommand{\emph}[1]{\textsf{\textbf{#1}}}
\usepackage{tikz}
\usetikzlibrary{calc,trees,positioning,arrows,fit,shapes,through, backgrounds}
\usetikzlibrary{patterns}

\usetikzlibrary{decorations.markings}
\usetikzlibrary{arrows}

\usepackage{pgfplots}

\usepackage{longtable}
\usepackage{tabularx}

\newcommand{\ds}{\displaystyle}
\newcommand{\ans}[1][1in]{\rule{#1}{.5pt}}

\newcommand{\points}[1]{(#1 points.)}		% Trying to be lazy.

\usepackage{array}
\newcolumntype{L}[1]{>{\raggedright\let\newline\\\arraybackslash\hspace{0pt}}m{#1}}
\newcolumntype{C}[1]{>{\centering\let\newline\\\arraybackslash\hspace{0pt}}m{#1}}
\newcolumntype{R}[1]{>{\raggedleft\let\newline\\\arraybackslash\hspace{0pt}}m{#1}}
\newcommand{\red}[1]{\textcolor{red}{#1}}

\newcommand{\be}{\begin{enumerate}}
\newcommand{\ee}{\end{enumerate}}

\begin{document}

\begin{enumerate}
\item \emph{Double Transposition Cipher:} Use two keywords to do a tabular transposition cipher twice.\\

\emph{Encryption Steps}

\be[i.]
\item Write the plaintext in \textbf{rows} using Keyword 1.
\item Rewrite columns with Keyword 1 in alphabetical order.
\item Read off the \emph{columns} to form the new ``plaintext''. 
\item Repeat this process with Keyword 2.\\
\ee


\item \emph{Example:} Encrypt the word \textbf{TRANSPOSITION} using double transposition with \begin{center} Keyword 1: {\tt KEY} and Keyword 2: {\tt SHORT}. \end{center}

\vfill

\item How is this different from the tabular transposition with keyword encryption method we used earlier?\\

\vspace{.5in}

\newpage


\item \emph{Decryption:}

\be[i.]
\item Make four grid headers: Keyword 2 (alphabetical order), Keyword 2 (natural order), Keyword 1 (alphabetical order), Keyword 1 (natural order).
\item Count the number of characters in the ciphertext and determine grid dimensions.
\item Starting with Keyword 2 (alphabetical) grid, fill  \emph{columns}.
\item Copy columns into the grid for Keyword 2 (natural order).
\item Read off the {\bf rows} to find new ciphertext.
\item Repeat (iii) with Keyword 1 (alphabetical).
\ee


\item {Decrypt} the ciphertext from problem 2. (We know it should decrypt as TRANSPOSITION!)


\vfill


\item {Decrypt} the ciphertext \so{ETETS PIHXC R} assuming it was encrypted using the scheme above.

\vfill



\end{enumerate}
\end{document}

%-------------------------------------------------------------------------------------------------------------------------------------------------------------------------------------------------------------------

%%% Local Variables:
%%% mode: latex
%%% TeX-master: t
%%% End:
